\subsection{FFLAGS and FSTATUS Registers}
FFLAGS holds accrued IEEE-754 (:term:base.terms.floating_point_exception_flag|plural:). FSTATUS holds the matching exception-condition enables, rounding mode, and optional
non-default floating-point modes. Both registers are present only with the floating-point extension and reset to zero.

\begin{manuallistedformatdiagram}{FFLAGS Register Format}
\manualformatrow{FFLAGS[15:0]}{%
\manualformatreserved{(:term:base.terms.reserved:)}{11}
\manualformatfield{I}{1}
\manualformatfield{Z}{1}
\manualformatfield{O}{1}
\manualformatfield{U}{1}
\manualformatfield{X}{1}
}
\end{manuallistedformatdiagram}

\begin{manuallistedformatdiagram}{FSTATUS Register Format}
\manualformatrow{FSTATUS[15:0]}{%
\manualformatreserved{(:term:base.terms.reserved:)}{6}
\manualformatfield{D}{1}
\manualformatfield{A}{1}
\manualformatfield{F}{1}
\manualformatfield{R}{2}
\manualformatfield{I}{1}
\manualformatfield{Z}{1}
\manualformatfield{O}{1}
\manualformatfield{U}{1}
\manualformatfield{X}{1}
}
\end{manuallistedformatdiagram}
{\small
In the diagrams, \texttt{I}, \texttt{Z}, \texttt{O}, \texttt{U}, and \texttt{X} denote the
\texttt{NV}, \texttt{DZ}, \texttt{OF}, \texttt{UF}, and \texttt{NX} flag or enable bits.
In FSTATUS, \texttt{D}, \texttt{A}, \texttt{F}, and \texttt{R} denote
\texttt{DN}, \texttt{DAZ}, \texttt{FTZ}, and \texttt{RM}.
\par}

\manualtablecaption{Floating-Point State Fields}
\begin{manualtabular}{@{}>{\raggedright\arraybackslash}p{0.75in}>{\raggedright\arraybackslash}p{0.65in}>{\raggedright\arraybackslash}p{3.95in}@{}}
\toprule
\rowcolor{ManualHeaderFill}
\textbf{Field} & \textbf{Bits} & \textbf{Meaning}\\
\midrule
\texttt{NV..NX} & \texttt{4..0} & accrued invalid, divide-by-zero, overflow, underflow, and inexact flags\\
\texttt{NV\_EN..NX\_EN} & \texttt{4..0} & floating-point exception-condition enables corresponding to FFLAGS.NV through FFLAGS.NX\\
\texttt{RM} & \texttt{6..5} & 0 nearest-even; 1 toward zero; 2 toward negative; 3 toward positive\\
\texttt{FTZ} & \texttt{7} & flush tiny results to signed zero\\
\texttt{DAZ} & \texttt{8} & treat subnormal inputs as signed zero\\
\texttt{DN} & \texttt{9} & produce the architectural default NaN\\
(:term:base.terms.reserved:) & \texttt{15..10} & (:term:base.terms.must_be_zero:)\\
\bottomrule
\end{manualtabular}
